\scanpage{100}
\margin{\textcircled{26}}
Formulaire faisant intervenir $\overset{\mathbb L}{\otimes}$ et $D^-$

\textcircled{1} Généralités valables sur tout topos annelé

Module \uline{plat} $F$ : $F\otimes\cdot$ est exact
\[
F\text{ plat}\Longleftrightarrow\text{ Pour tout injectif }I,\quad\mathcal Hom(F,I)\text{ est injectif}
\]
\margin{\uline{NB} Dans le cas actuel, il suffit de vérifier sur les fibres.}
Tout Module $F$ \deleted{admet des \unsure{ré}} est quotient d'un plat, donc admet des résolutions plates : on prend les $A_U$ ($U\in\operatorname{Ob}X_{\mathrm{ét}}$). Il s'ensuit \inserted{\unsure{formellement}} que pour $\forall\,L_\bullet\in\operatorname{Ob}D^-(X)$, existe \unsure{résoln} $L'_\bullet\unsure{-|\to}L_\bullet$, avec $L'_\bullet$ plat en tous degrés, qui est un quasi-isomorphisme.
De ceci résulte formellement que $D^-(X)$ est équivalent : la catégorie obtenue en prenant \inserted{la catégorie $K^-(X)^{\mathrm{plat}}$} des complexes $L_\bullet$ bornés sup., plats \deleted{en} tous degrés, \deleted{\unsure{et}} en prenant les hom. à homotopie près, puis prenant un quotient par la \deleted{relation} \inserted{famille} des quasi-isomorphismes.

Ceci va nous permettre de prendre les ``foncteurs dérivés gauches'' du foncteur $\otimes$, en prenant dans $C^-(X)^{\mathrm{plat}}$ le $\otimes$ ordinaire; \inserted{on notera que ça passe au quotient par quasi-isomorphisme.}
Donc si $L_\bullet,M_\bullet$ sont deux complexes ordinaires $\in\operatorname{Ob}D^-(X)$, on \deleted{définit} calcule $L_\bullet\overset{L}{\otimes}M_\bullet$ en prenant des résolutions plates $L'_\bullet\to L_\bullet$, $M'_\bullet\unsure{-|\to}M_\bullet$ et prenant $L'_\bullet\otimes M'_\bullet$.
\mspage{99}
